% --- Archon physics formalization source begin ---
% archon:physics
% archon:covers PhyXMiniProblems/problem_phyx_mini_0427.lean
% archon:source-report reports/phyx_mini/problem_phyx_mini_0427.source.json

\chapter{Physics problem phyx\_mini\_0427}
\label{ch:PhyXMiniProblems_problem_phyx_mini_0427}

\paragraph{Problem source.}
\#\# Physical scenario

Figure shows the cycle for a heat engine that uses a gas having \$\textbackslash{}gamma = 1.25\$. The initial temperature is \$T\_1 = 300 \textbackslash{}, \textbackslash{}text\{K\}\$, and this engine operates at 20 cycles per second.

\#\# Question

What is the engine's thermal efficiency?

\#\# Answer choices

- A: A: 0.668

- B: B: 0.15

- C: C: 0.548

- D: D: 0.0901

\#\# Figure caption (auxiliary; use the image as primary evidence)

The image is a pressure-volume (p-V) diagram commonly used in thermodynamics. It includes the following elements:

1. **Axes:**
   - The vertical axis represents pressure (p) in atmospheres (atm).
   - The horizontal axis represents volume (V) in cubic centimeters (cm³).

2. **Points and Numbers:** 
   - The graph has three numbered points labeled as 1, 2, and 3. These points are connected, representing different states in a thermodynamic process.
   - **Point 2** is at the top left.
   - **Point 3** is at the top right.
   - **Point 1** is at the bottom right.

3. **Arrows:**
   - An arrow connects point 2 to point 3, indicating a process at constant pressure.
   - Another arrow points downward from point 3 to point 1, indicating a process at constant volume.
   - A curved arrow connects point 1 back to point 2, labeled "Isotherm," indicating an isothermal process.

4. **Text and Labels:**
   - "p (atm)" is labeled on the pressure axis, with a specific label "p\_max" indicating maximum pressure.
   - "V (cm³)" is labeled on the volume axis.
   - "Isotherm" is labeled near the curve connecting points 1 and 2, indicating that this path occurs at a constant temperature.

The diagram illustrates a typical cyclical process involving isochoric, isobaric, and isothermal transformations.

\#\# Dataset metadata

Laws of Thermodynamics; Physical Model Grounding Reasoning, Multi-Formula Reasoning

\paragraph{Recorded answer/context.}
D: D: 0.0901

\paragraph{Figure/image path.}
/root/proposal\_for\_physic/science-mango/phyx\_mini\_run/phyx\_data/test\_image/427.png

\paragraph{Formalization target.}
create a compiling Lean file with sorry bodies at `PhyXMiniProblems/problem\_phyx\_mini\_0427.lean`. The Lean declarations must preserve the physical quantities, dimensions or dimensional roles, figure labels, governing-law hypotheses, and final relation expressed by this problem.
Use Mathlib/Physlib names found through LeanExplore where available. If a physics API is missing, introduce faithful local abstractions rather than scalar placeholder aliases.

\begin{theorem}[Physics formalization target]
\label{thm:physics:phyx_mini_0427:target}
\lean{PhyXMiniProblems.ProblemPhyXMini0427.problem_phyx_mini_0427}
\uses{def:physics:phyx-mini-0427:phyxminiproblems-problemphyxmini0427-idealgasheatenginecycle, def:physics:phyx-mini-0427:phyxminiproblems-problemphyxmini0427-matchesproblemstatement, def:physics:phyx-mini-0427:phyxminiproblems-problemphyxmini0427-matchesprimarypressurevolumefigure, def:physics:phyx-mini-0427:phyxminiproblems-problemphyxmini0427-hasphysicalthermodynamicparameters, def:physics:phyx-mini-0427:phyxminiproblems-problemphyxmini0427-obeyscaloricallyperfectidealgaslaws, def:physics:phyx-mini-0427:phyxminiproblems-problemphyxmini0427-thermalefficiency, def:physics:phyx-mini-0427:phyxminiproblems-problemphyxmini0427-recordedanswerchoice, def:physics:phyx-mini-0427:phyxminiproblems-problemphyxmini0427-matchestonearesttenthousandth, def:physics:phyx-mini-0427:phyxminiproblems-problemphyxmini0427-isuniquematchingdisplayedefficiency}
The assigned autoformalize agent should translate this physics problem into Lean declarations in the covered file, with theorem and lemma proof bodies written as `by sorry`.
\end{theorem}
\begin{proof}
This is an autoformalization task, not a proof task. Produce statements that can later be proved by the physics prover without weakening the physical model.
\end{proof}

% --- Archon declaration topology begin ---
\section{Declaration topology}
\begin{definition}[Volume Quantity]
  \label{def:physics:phyx-mini-0427:phyxminiproblems-problemphyxmini0427-volumequantity}
  \lean{PhyXMiniProblems.ProblemPhyXMini0427.VolumeQuantity}
  A nonnegative physical volume carrying dimension L³
\end{definition}
\begin{proof}
  This is the declared model definition of Volume Quantity.
\end{proof}
\begin{definition}[Cycle Rate Quantity]
  \label{def:physics:phyx-mini-0427:phyxminiproblems-problemphyxmini0427-cycleratequantity}
  \lean{PhyXMiniProblems.ProblemPhyXMini0427.CycleRateQuantity}
  A nonnegative cycle rate carrying inverse-time dimension
\end{definition}
\begin{proof}
  This is the declared model definition of Cycle Rate Quantity.
\end{proof}
\begin{definition}[pressure In Pascals]
  \label{def:physics:phyx-mini-0427:phyxminiproblems-problemphyxmini0427-pressureinpascals}
  \lean{PhyXMiniProblems.ProblemPhyXMini0427.pressureInPascals}
  Read a physical pressure in coherent SI pascals
\end{definition}
\begin{proof}
  This is the declared model definition of pressure In Pascals.
\end{proof}
\begin{definition}[pressure In Atmospheres]
  \label{def:physics:phyx-mini-0427:phyxminiproblems-problemphyxmini0427-pressureinatmospheres}
  \lean{PhyXMiniProblems.ProblemPhyXMini0427.pressureInAtmospheres}
  \uses{def:physics:phyx-mini-0427:phyxminiproblems-problemphyxmini0427-pressureinpascals}
  Read a pressure in the atmospheres printed on the vertical axis
\end{definition}
\begin{proof}
  This is the declared model definition of pressure In Atmospheres.
\end{proof}
\begin{definition}[volume In Cubic Metres]
  \label{def:physics:phyx-mini-0427:phyxminiproblems-problemphyxmini0427-volumeincubicmetres}
  \lean{PhyXMiniProblems.ProblemPhyXMini0427.volumeInCubicMetres}
  \uses{def:physics:phyx-mini-0427:phyxminiproblems-problemphyxmini0427-volumequantity}
  Read a physical volume in coherent SI cubic metres
\end{definition}
\begin{proof}
  This is the declared model definition of volume In Cubic Metres.
\end{proof}
\begin{definition}[volume In Cubic Centimetres]
  \label{def:physics:phyx-mini-0427:phyxminiproblems-problemphyxmini0427-volumeincubiccentimetres}
  \lean{PhyXMiniProblems.ProblemPhyXMini0427.volumeInCubicCentimetres}
  \uses{def:physics:phyx-mini-0427:phyxminiproblems-problemphyxmini0427-volumequantity, def:physics:phyx-mini-0427:phyxminiproblems-problemphyxmini0427-volumeincubicmetres}
  Read a volume in the cubic centimetres printed on the horizontal axis
\end{definition}
\begin{proof}
  This is the declared model definition of volume In Cubic Centimetres.
\end{proof}
\begin{definition}[energy In Joules]
  \label{def:physics:phyx-mini-0427:phyxminiproblems-problemphyxmini0427-energyinjoules}
  \lean{PhyXMiniProblems.ProblemPhyXMini0427.energyInJoules}
  Read signed heat, work, or internal-energy change in SI joules
\end{definition}
\begin{proof}
  This is the declared model definition of energy In Joules.
\end{proof}
\begin{definition}[temperature In Kelvin]
  \label{def:physics:phyx-mini-0427:phyxminiproblems-problemphyxmini0427-temperatureinkelvin}
  \lean{PhyXMiniProblems.ProblemPhyXMini0427.temperatureInKelvin}
  Read an absolute temperature in kelvins from its stated storage unit
\end{definition}
\begin{proof}
  This is the declared model definition of temperature In Kelvin.
\end{proof}
\begin{definition}[cycle Rate In Hertz]
  \label{def:physics:phyx-mini-0427:phyxminiproblems-problemphyxmini0427-cyclerateinhertz}
  \lean{PhyXMiniProblems.ProblemPhyXMini0427.cycleRateInHertz}
  \uses{def:physics:phyx-mini-0427:phyxminiproblems-problemphyxmini0427-cycleratequantity}
  Read a physical cycle rate in cycles per SI second (hertz)
\end{definition}
\begin{proof}
  This is the declared model definition of cycle Rate In Hertz.
\end{proof}
\begin{definition}[Gas Model]
  \label{def:physics:phyx-mini-0427:phyxminiproblems-problemphyxmini0427-gasmodel}
  \lean{PhyXMiniProblems.ProblemPhyXMini0427.GasModel}
  The working-substance model used for the gas
\end{definition}
\begin{proof}
  This is the declared model definition of Gas Model.
\end{proof}
\begin{definition}[Thermodynamic Device Role]
  \label{def:physics:phyx-mini-0427:phyxminiproblems-problemphyxmini0427-thermodynamicdevicerole}
  \lean{PhyXMiniProblems.ProblemPhyXMini0427.ThermodynamicDeviceRole}
  The thermodynamic role played by the cyclic device
\end{definition}
\begin{proof}
  This is the declared model definition of Thermodynamic Device Role.
\end{proof}
\begin{definition}[State Label]
  \label{def:physics:phyx-mini-0427:phyxminiproblems-problemphyxmini0427-statelabel}
  \lean{PhyXMiniProblems.ProblemPhyXMini0427.StateLabel}
  The three numbered states printed in the supplied diagram
\end{definition}
\begin{proof}
  This is the declared model definition of State Label.
\end{proof}
\begin{definition}[Cycle Leg]
  \label{def:physics:phyx-mini-0427:phyxminiproblems-problemphyxmini0427-cycleleg}
  \lean{PhyXMiniProblems.ProblemPhyXMini0427.CycleLeg}
  The three directed arrows making up one traversal of the cycle
\end{definition}
\begin{proof}
  This is the declared model definition of Cycle Leg.
\end{proof}
\begin{definition}[initial State]
  \label{def:physics:phyx-mini-0427:phyxminiproblems-problemphyxmini0427-cycleleg-initialstate}
  \lean{PhyXMiniProblems.ProblemPhyXMini0427.CycleLeg.initialState}
  \uses{def:physics:phyx-mini-0427:phyxminiproblems-problemphyxmini0427-statelabel, def:physics:phyx-mini-0427:phyxminiproblems-problemphyxmini0427-cycleleg}
  Initial state of each directed process leg
\end{definition}
\begin{proof}
  This is the declared model definition of initial State.
\end{proof}
\begin{definition}[final State]
  \label{def:physics:phyx-mini-0427:phyxminiproblems-problemphyxmini0427-cycleleg-finalstate}
  \lean{PhyXMiniProblems.ProblemPhyXMini0427.CycleLeg.finalState}
  \uses{def:physics:phyx-mini-0427:phyxminiproblems-problemphyxmini0427-statelabel, def:physics:phyx-mini-0427:phyxminiproblems-problemphyxmini0427-cycleleg}
  Final state of each directed process leg
\end{definition}
\begin{proof}
  This is the declared model definition of final State.
\end{proof}
\begin{definition}[Process Kind]
  \label{def:physics:phyx-mini-0427:phyxminiproblems-problemphyxmini0427-processkind}
  \lean{PhyXMiniProblems.ProblemPhyXMini0427.ProcessKind}
  Thermodynamic constraint holding on a process leg
\end{definition}
\begin{proof}
  This is the declared model definition of Process Kind.
\end{proof}
\begin{definition}[Path Geometry]
  \label{def:physics:phyx-mini-0427:phyxminiproblems-problemphyxmini0427-pathgeometry}
  \lean{PhyXMiniProblems.ProblemPhyXMini0427.PathGeometry}
  Geometric appearance of a process in the pressure--volume plane
\end{definition}
\begin{proof}
  This is the declared model definition of Path Geometry.
\end{proof}
\begin{definition}[Axis Quantity]
  \label{def:physics:phyx-mini-0427:phyxminiproblems-problemphyxmini0427-axisquantity}
  \lean{PhyXMiniProblems.ProblemPhyXMini0427.AxisQuantity}
  Physical quantity assigned to one of the two plot axes
\end{definition}
\begin{proof}
  This is the declared model definition of Axis Quantity.
\end{proof}
\begin{definition}[Axis Unit]
  \label{def:physics:phyx-mini-0427:phyxminiproblems-problemphyxmini0427-axisunit}
  \lean{PhyXMiniProblems.ProblemPhyXMini0427.AxisUnit}
  Unit explicitly printed beside an axis
\end{definition}
\begin{proof}
  This is the declared model definition of Axis Unit.
\end{proof}
\begin{definition}[Figure Label]
  \label{def:physics:phyx-mini-0427:phyxminiproblems-problemphyxmini0427-figurelabel}
  \lean{PhyXMiniProblems.ProblemPhyXMini0427.FigureLabel}
  Text and numerical labels visible in the primary bitmap
\end{definition}
\begin{proof}
  This is the declared model definition of Figure Label.
\end{proof}
\begin{definition}[Thermodynamic State]
  \label{def:physics:phyx-mini-0427:phyxminiproblems-problemphyxmini0427-thermodynamicstate}
  \lean{PhyXMiniProblems.ProblemPhyXMini0427.ThermodynamicState}
  \uses{def:physics:phyx-mini-0427:phyxminiproblems-problemphyxmini0427-volumequantity}
  A physical equilibrium state of the working gas
\end{definition}
\begin{proof}
  This is the declared model definition of Thermodynamic State.
\end{proof}
\begin{definition}[Pressure Volume Figure]
  \label{def:physics:phyx-mini-0427:phyxminiproblems-problemphyxmini0427-pressurevolumefigure}
  \lean{PhyXMiniProblems.ProblemPhyXMini0427.PressureVolumeFigure}
  \uses{def:physics:phyx-mini-0427:phyxminiproblems-problemphyxmini0427-statelabel, def:physics:phyx-mini-0427:phyxminiproblems-problemphyxmini0427-cycleleg, def:physics:phyx-mini-0427:phyxminiproblems-problemphyxmini0427-pathgeometry, def:physics:phyx-mini-0427:phyxminiproblems-problemphyxmini0427-axisquantity, def:physics:phyx-mini-0427:phyxminiproblems-problemphyxmini0427-axisunit, def:physics:phyx-mini-0427:phyxminiproblems-problemphyxmini0427-figurelabel}
  Qualitative plot data kept separate from the physical state values. In particular, this structure records axes, labels, vertices, path shapes, and arrow directions, but no efficiency or answer choice
\end{definition}
\begin{proof}
  This is the declared model definition of Pressure Volume Figure.
\end{proof}
\begin{definition}[Ideal Gas Heat Engine Cycle]
  \label{def:physics:phyx-mini-0427:phyxminiproblems-problemphyxmini0427-idealgasheatenginecycle}
  \lean{PhyXMiniProblems.ProblemPhyXMini0427.IdealGasHeatEngineCycle}
  \uses{def:physics:phyx-mini-0427:phyxminiproblems-problemphyxmini0427-cycleratequantity, def:physics:phyx-mini-0427:phyxminiproblems-problemphyxmini0427-gasmodel, def:physics:phyx-mini-0427:phyxminiproblems-problemphyxmini0427-thermodynamicdevicerole, def:physics:phyx-mini-0427:phyxminiproblems-problemphyxmini0427-statelabel, def:physics:phyx-mini-0427:phyxminiproblems-problemphyxmini0427-cycleleg, def:physics:phyx-mini-0427:phyxminiproblems-problemphyxmini0427-processkind, def:physics:phyx-mini-0427:phyxminiproblems-problemphyxmini0427-thermodynamicstate, def:physics:phyx-mini-0427:phyxminiproblems-problemphyxmini0427-pressurevolumefigure}
  Independent physical observables for one fixed sample of gas. Work is done by the gas, heat is transferred into the gas, and internal-energy change is U\_final - U\_initial on each leg. The requested efficiency is deliberately not stored as a field
\end{definition}
\begin{proof}
  This is the declared model definition of Ideal Gas Heat Engine Cycle.
\end{proof}
\begin{definition}[Matches Problem Statement]
  \label{def:physics:phyx-mini-0427:phyxminiproblems-problemphyxmini0427-matchesproblemstatement}
  \lean{PhyXMiniProblems.ProblemPhyXMini0427.MatchesProblemStatement}
  \uses{def:physics:phyx-mini-0427:phyxminiproblems-problemphyxmini0427-temperatureinkelvin, def:physics:phyx-mini-0427:phyxminiproblems-problemphyxmini0427-cyclerateinhertz, def:physics:phyx-mini-0427:phyxminiproblems-problemphyxmini0427-idealgasheatenginecycle}
  Numerical data stated in the prose. The cycle rate is retained even though thermal efficiency is a per-cycle ratio and hence independent of that rate. No efficiency or answer value occurs here
\end{definition}
\begin{proof}
  This is the declared model definition of Matches Problem Statement.
\end{proof}
\begin{definition}[Matches Primary Pressure Volume Figure]
  \label{def:physics:phyx-mini-0427:phyxminiproblems-problemphyxmini0427-matchesprimarypressurevolumefigure}
  \lean{PhyXMiniProblems.ProblemPhyXMini0427.MatchesPrimaryPressureVolumeFigure}
  \uses{def:physics:phyx-mini-0427:phyxminiproblems-problemphyxmini0427-pressureinatmospheres, def:physics:phyx-mini-0427:phyxminiproblems-problemphyxmini0427-volumeincubiccentimetres, def:physics:phyx-mini-0427:phyxminiproblems-problemphyxmini0427-statelabel, def:physics:phyx-mini-0427:phyxminiproblems-problemphyxmini0427-cycleleg, def:physics:phyx-mini-0427:phyxminiproblems-problemphyxmini0427-figurelabel, def:physics:phyx-mini-0427:phyxminiproblems-problemphyxmini0427-idealgasheatenginecycle}
  Exact transcription of image 427.png. The bitmap fixes state 1 at (600 cm³, 1 atm), state 2 at volume 200 cm³, state 3 at volume 600 cm³, and states 2 and 3 on the common line labelled p\_max. The numerical value of p\_max is not assumed; it must follow from the isothermal ideal-gas leg
\end{definition}
\begin{proof}
  This is the declared model definition of Matches Primary Pressure Volume Figure.
\end{proof}
\begin{definition}[Has Physical Thermodynamic Parameters]
  \label{def:physics:phyx-mini-0427:phyxminiproblems-problemphyxmini0427-hasphysicalthermodynamicparameters}
  \lean{PhyXMiniProblems.ProblemPhyXMini0427.HasPhysicalThermodynamicParameters}
  \uses{def:physics:phyx-mini-0427:phyxminiproblems-problemphyxmini0427-pressureinpascals, def:physics:phyx-mini-0427:phyxminiproblems-problemphyxmini0427-volumeincubicmetres, def:physics:phyx-mini-0427:phyxminiproblems-problemphyxmini0427-temperatureinkelvin, def:physics:phyx-mini-0427:phyxminiproblems-problemphyxmini0427-cyclerateinhertz, def:physics:phyx-mini-0427:phyxminiproblems-problemphyxmini0427-idealgasheatenginecycle}
  Positivity and nondegeneracy of the physical parameters
\end{definition}
\begin{proof}
  This is the declared model definition of Has Physical Thermodynamic Parameters.
\end{proof}
\begin{definition}[Obeys Calorically Perfect Ideal Gas Laws]
  \label{def:physics:phyx-mini-0427:phyxminiproblems-problemphyxmini0427-obeyscaloricallyperfectidealgaslaws}
  \lean{PhyXMiniProblems.ProblemPhyXMini0427.ObeysCaloricallyPerfectIdealGasLaws}
  \uses{def:physics:phyx-mini-0427:phyxminiproblems-problemphyxmini0427-pressureinpascals, def:physics:phyx-mini-0427:phyxminiproblems-problemphyxmini0427-volumeincubicmetres, def:physics:phyx-mini-0427:phyxminiproblems-problemphyxmini0427-energyinjoules, def:physics:phyx-mini-0427:phyxminiproblems-problemphyxmini0427-temperatureinkelvin, def:physics:phyx-mini-0427:phyxminiproblems-problemphyxmini0427-statelabel, def:physics:phyx-mini-0427:phyxminiproblems-problemphyxmini0427-cycleleg, def:physics:phyx-mini-0427:phyxminiproblems-problemphyxmini0427-idealgasheatenginecycle}
  Macroscopic governing laws for a fixed sample of calorically perfect ideal gas with constant heat-capacity ratio γ. the cross-multiplied fixed-sample ideal-gas law expresses pV/T = const; ΔU = (p\_f V\_f - p\_i V\_i)/(γ - 1); the first law uses Q\_into = ΔU + W\_by; reversible isothermal work is p\_i V\_i log(V\_f/V\_i); isobaric work is p(V\_f-V\_i); an isochoric leg has zero boundary work. These are general laws.
\end{definition}
\begin{proof}
  This is the declared model definition of Obeys Calorically Perfect Ideal Gas Laws.
\end{proof}
\begin{definition}[initial Pressure Volume Energy In Joules]
  \label{def:physics:phyx-mini-0427:phyxminiproblems-problemphyxmini0427-initialpressurevolumeenergyinjoules}
  \lean{PhyXMiniProblems.ProblemPhyXMini0427.initialPressureVolumeEnergyInJoules}
  \uses{def:physics:phyx-mini-0427:phyxminiproblems-problemphyxmini0427-pressureinpascals, def:physics:phyx-mini-0427:phyxminiproblems-problemphyxmini0427-volumeincubicmetres, def:physics:phyx-mini-0427:phyxminiproblems-problemphyxmini0427-idealgasheatenginecycle}
  The p₁V₁ energy scale of this cycle, expressed in joules
\end{definition}
\begin{proof}
  This is the declared model definition of initial Pressure Volume Energy In Joules.
\end{proof}
\begin{definition}[net Work Done By Gas In Joules]
  \label{def:physics:phyx-mini-0427:phyxminiproblems-problemphyxmini0427-networkdonebygasinjoules}
  \lean{PhyXMiniProblems.ProblemPhyXMini0427.netWorkDoneByGasInJoules}
  \uses{def:physics:phyx-mini-0427:phyxminiproblems-problemphyxmini0427-energyinjoules, def:physics:phyx-mini-0427:phyxminiproblems-problemphyxmini0427-idealgasheatenginecycle}
  Net work done by the gas in one clockwise traversal of the cycle
\end{definition}
\begin{proof}
  This is the declared model definition of net Work Done By Gas In Joules.
\end{proof}
\begin{definition}[total Heat Absorbed In Joules]
  \label{def:physics:phyx-mini-0427:phyxminiproblems-problemphyxmini0427-totalheatabsorbedinjoules}
  \lean{PhyXMiniProblems.ProblemPhyXMini0427.totalHeatAbsorbedInJoules}
  \uses{def:physics:phyx-mini-0427:phyxminiproblems-problemphyxmini0427-energyinjoules, def:physics:phyx-mini-0427:phyxminiproblems-problemphyxmini0427-idealgasheatenginecycle}
  Total heat absorbed per cycle, summing the positive parts of signed heat
\end{definition}
\begin{proof}
  This is the declared model definition of total Heat Absorbed In Joules.
\end{proof}
\begin{definition}[thermal Efficiency]
  \label{def:physics:phyx-mini-0427:phyxminiproblems-problemphyxmini0427-thermalefficiency}
  \lean{PhyXMiniProblems.ProblemPhyXMini0427.thermalEfficiency}
  \uses{def:physics:phyx-mini-0427:phyxminiproblems-problemphyxmini0427-idealgasheatenginecycle, def:physics:phyx-mini-0427:phyxminiproblems-problemphyxmini0427-networkdonebygasinjoules, def:physics:phyx-mini-0427:phyxminiproblems-problemphyxmini0427-totalheatabsorbedinjoules}
  Dimensionless thermal efficiency W\_net/Q\_absorbed for one cycle
\end{definition}
\begin{proof}
  This is the declared model definition of thermal Efficiency.
\end{proof}
\begin{definition}[Answer Choice]
  \label{def:physics:phyx-mini-0427:phyxminiproblems-problemphyxmini0427-answerchoice}
  \lean{PhyXMiniProblems.ProblemPhyXMini0427.AnswerChoice}
  The four answer labels printed in the source problem
\end{definition}
\begin{proof}
  This is the declared model definition of Answer Choice.
\end{proof}
\begin{definition}[displayed Efficiency]
  \label{def:physics:phyx-mini-0427:phyxminiproblems-problemphyxmini0427-displayedefficiency}
  \lean{PhyXMiniProblems.ProblemPhyXMini0427.displayedEfficiency}
  \uses{def:physics:phyx-mini-0427:phyxminiproblems-problemphyxmini0427-answerchoice}
  Dimensionless decimal printed beside each answer label
\end{definition}
\begin{proof}
  This is the declared model definition of displayed Efficiency.
\end{proof}
\begin{definition}[recorded Answer Choice]
  \label{def:physics:phyx-mini-0427:phyxminiproblems-problemphyxmini0427-recordedanswerchoice}
  \lean{PhyXMiniProblems.ProblemPhyXMini0427.recordedAnswerChoice}
  \uses{def:physics:phyx-mini-0427:phyxminiproblems-problemphyxmini0427-answerchoice}
  Answer label recorded in the dataset metadata
\end{definition}
\begin{proof}
  This is the declared model definition of recorded Answer Choice.
\end{proof}
\begin{definition}[Matches To Nearest Ten Thousandth]
  \label{def:physics:phyx-mini-0427:phyxminiproblems-problemphyxmini0427-matchestonearesttenthousandth}
  \lean{PhyXMiniProblems.ProblemPhyXMini0427.MatchesToNearestTenThousandth}
  \uses{def:physics:phyx-mini-0427:phyxminiproblems-problemphyxmini0427-idealgasheatenginecycle, def:physics:phyx-mini-0427:phyxminiproblems-problemphyxmini0427-thermalefficiency, def:physics:phyx-mini-0427:phyxminiproblems-problemphyxmini0427-answerchoice, def:physics:phyx-mini-0427:phyxminiproblems-problemphyxmini0427-displayedefficiency}
  Agreement with a decimal displayed to four places
\end{definition}
\begin{proof}
  This is the declared model definition of Matches To Nearest Ten Thousandth.
\end{proof}
\begin{definition}[Is Unique Matching Displayed Efficiency]
  \label{def:physics:phyx-mini-0427:phyxminiproblems-problemphyxmini0427-isuniquematchingdisplayedefficiency}
  \lean{PhyXMiniProblems.ProblemPhyXMini0427.IsUniqueMatchingDisplayedEfficiency}
  \uses{def:physics:phyx-mini-0427:phyxminiproblems-problemphyxmini0427-idealgasheatenginecycle, def:physics:phyx-mini-0427:phyxminiproblems-problemphyxmini0427-answerchoice, def:physics:phyx-mini-0427:phyxminiproblems-problemphyxmini0427-matchestonearesttenthousandth}
  The specified choice is the unique four-decimal display match
\end{definition}
\begin{proof}
  This is the declared model definition of Is Unique Matching Displayed Efficiency.
\end{proof}
\begin{lemma}[cycle State And Energy Readouts]
  \label{lem:physics:phyx-mini-0427:phyxminiproblems-problemphyxmini0427-cyclestateandenergyreadouts}
  \lean{PhyXMiniProblems.ProblemPhyXMini0427.cycleStateAndEnergyReadouts}
  \uses{def:physics:phyx-mini-0427:phyxminiproblems-problemphyxmini0427-pressureinatmospheres, def:physics:phyx-mini-0427:phyxminiproblems-problemphyxmini0427-energyinjoules, def:physics:phyx-mini-0427:phyxminiproblems-problemphyxmini0427-temperatureinkelvin, def:physics:phyx-mini-0427:phyxminiproblems-problemphyxmini0427-idealgasheatenginecycle, def:physics:phyx-mini-0427:phyxminiproblems-problemphyxmini0427-matchesproblemstatement, def:physics:phyx-mini-0427:phyxminiproblems-problemphyxmini0427-matchesprimarypressurevolumefigure, def:physics:phyx-mini-0427:phyxminiproblems-problemphyxmini0427-hasphysicalthermodynamicparameters, def:physics:phyx-mini-0427:phyxminiproblems-problemphyxmini0427-obeyscaloricallyperfectidealgaslaws, def:physics:phyx-mini-0427:phyxminiproblems-problemphyxmini0427-initialpressurevolumeenergyinjoules}
  The figure and generic laws imply p\_max = 3 atm, T₂ = 300 K, and T₃ = 900 K. Writing E₁ = p₁V₁, the leg heats are -E₁ log 3, 10E₁, and -8E₁, while the works are -E₁ log 3, 2E₁, and 0. These are derived conclusions, not premise fields
\end{lemma}
\begin{proof}
  The conclusion follows from the governing relations and figure readouts named in the statement of cycle State And Energy Readouts.
\end{proof}
% --- Archon declaration topology end ---
% --- Archon physics formalization source end ---
